% =====================================================================
%  chapters/minggu-13.tex
%  Minggu 13 -- Pengenalan React: Komponen, JSX & Props
% =====================================================================

\chapter{Pengenalan React: Komponen, JSX \& Props}
\label{chap:mjs-13}

% ---------------------------------------------------------------------
% 1. CAPAIAN PEMBELAJARAN
% ---------------------------------------------------------------------
\begin{capaian}
Setelah menyelesaikan bab ini, mahasiswa diharapkan mampu:
\begin{itemize}
  \item menjelaskan mengapa React digunakan dan bagaimana posisinya dalam pengembangan web;
  \item membaca dan memahami kode React sederhana (komponen, JSX, \emph{props});
  \item membedakan JSX dari HTML biasa;
  \item mengidentifikasi struktur komponen, termasuk komponen bersarang dan penggunaan \texttt{map()} untuk daftar.
\end{itemize}
\end{capaian}

% ---------------------------------------------------------------------
% 2. TUJUAN PRAKTIK
% ---------------------------------------------------------------------
\begin{tujuanpraktik}
Pada sesi praktik ini, mahasiswa akan:
\begin{itemize}
  \item membaca kode React secara \emph{visual} tanpa menjalankannya;
  \item mengonstruksi file HTML berisi potongan kode React dengan \emph{syntax highlighting};
  \item berlatih menjawab pertanyaan pemahaman tentang kode React.
\end{itemize}
\end{tujuanpraktik}

\begin{catatan}
Modul ini \textbf{hanya pengenalan}. Kita \textbf{TIDAK} akan menginstal React, menjalankan \texttt{create-react-app} atau Vite, maupun menggunakan \emph{build tools}. Tujuannya: membaca, memahami, dan ``merasakan'' kode React agar konsepnya tidak asing saat belajar di semester lanjutan.
\end{catatan}

% =====================================================================
\section{Mengapa React?}
\label{sec:13-mengapa-react}
% =====================================================================

\textbf{React} adalah \emph{JavaScript library} untuk membangun \textbf{antarmuka (UI) berbasis komponen}---artinya, Anda memecah halaman web menjadi potongan-potongan kecil yang bisa digunakan ulang.

\subsection{Perbandingan: Tanpa React vs Dengan React}

\textbf{Tanpa React (vanilla JavaScript):}

\begin{lstlisting}[language=JavaScript, caption={Manipulasi DOM manual tanpa React}, label={lst:13-vanilla}]
// Anda harus secara manual mencari elemen dan mengubahnya
const daftar = document.getElementById("daftar-mahasiswa");
const item = document.createElement("li");
item.textContent = "Budi";
daftar.appendChild(item);
// Jika data berubah -> Anda harus menulis ulang kode manipulasi DOM
\end{lstlisting}

\textbf{Dengan React:}

\begin{lstlisting}[language=JavaScript, caption={Menulis UI secara deklaratif dengan React}, label={lst:13-react-basic}]
// Anda cukup mendeskripsikan APA yang ingin ditampilkan
function DaftarMahasiswa({ mahasiswa }) {
    return (
        <ul>
            {mahasiswa.map((nama, index) => (
                <li key={index}>{nama}</li>
            ))}
        </ul>
    );
}

// React yang mengurus bagaimana update DOM saat data berubah
\end{lstlisting}

\subsection{Mengapa React Populer?}

\begin{table}[htbp]
\centering
\caption{Alasan popularitas React}
\label{tab:13-mengapa-react}
\begin{tabular}{ll}
\toprule
\textbf{Alasan} & \textbf{Penjelasan} \\
\midrule
\emph{Komponen} & UI dipecah jadi blok kecil, bisa digunakan ulang \\
Deklaratif & Tulis ``apa'' yang ingin ditampilkan, bukan ``bagaimana'' update DOM \\
Ekosistem besar & Banyak library pendukung (routing, \emph{state management}, UI kit) \\
\emph{Industry standard} & Banyak perusahaan teknologi menggunakan React \\
\bottomrule
\end{tabular}
\end{table}

% =====================================================================
\section{Membaca Kode React (Tanpa Menjalankan)}
\label{sec:13-membaca-kode}
% =====================================================================

Tujuan: melihat kode React secara utuh agar terbiasa dengan sintaksnya. Kita \textbf{tidak menjalankan} kode ini---kita membaca dan memahaminya.

\begin{langkah}
  \item Buat file baru \texttt{react-preview.html}.
  \item Masukkan kode HTML berikut (berisi empat contoh komponen React).
  \item Buka di peramban untuk melihat tampilan kode yang sudah diberi warna \emph{syntax highlighting}.
\end{langkah}

\begin{tangkapanlayar}
  {assets/images/minggu-13/react-preview-kode.png}
  {Halaman \emph{preview} kode React dengan \emph{syntax highlighting} yang rapi.}
  {Buka \texttt{react-preview.html} di peramban setelah menambahkan kode-kode berikut.}
\end{tangkapanlayar}

\subsection{Kode 1 --- Komponen Paling Sederhana}

\begin{lstlisting}[language=JavaScript, caption={Komponen ``Halo Dunia'' paling sederhana}, label={lst:13-komponen-sederhana}]
// Dalam React, UI ditulis sebagai FUNCTION
// yang mengembalikan JSX
function Sapa() {
    return (
        <div className="container">
            <h1>Halo, Dunia React!</h1>
            <p>Ini adalah komponen React pertama saya.</p>
        </div>
    );
}
\end{lstlisting}

\begin{catatan}
JSX mirip HTML, tapi ada perbedaan kecil---misal: \texttt{class} menjadi \texttt{className}. Komponen harus diawali huruf KAPITAL (misal: \texttt{Sapa}, bukan \texttt{sapa}).
\end{catatan}

\subsection{Kode 2 --- Komponen dengan Props}

\begin{lstlisting}[language=JavaScript, caption={Komponen yang menerima data melalui \emph{props}}, label={lst:13-komponen-props}]
// Props = "parameter" untuk komponen React
function KartuProfil({ nama, jurusan, angkatan }) {
    return (
        <div className="bg-white rounded-lg shadow-md p-6">
            <h2 className="text-xl font-bold">{nama}</h2>
            <p className="text-gray-600">
                Jurusan: {jurusan}
            </p>
            <p className="text-gray-600">
                Angkatan: {angkatan}
            </p>
        </div>
    );
}

// Cara menggunakan (memanggil) komponen:
// <KartuProfil nama="Budi" jurusan="Informatika"
//              angkatan="2025" />
// <KartuProfil nama="Ani" jurusan="Informatika"
//              angkatan="2025" />
\end{lstlisting}

\begin{tip}
Fungsi \texttt{KartuProfil} menerima tiga \emph{props} (\texttt{nama}, \texttt{jurusan}, \texttt{angkatan}) dan menggunakannya di dalam JSX. Komponen yang sama bisa digunakan berulang kali dengan data berbeda---itulah kekuatan komponen!
\end{tip}

\subsection{Kode 3 --- Komponen dengan Daftar (Map)}

\begin{lstlisting}[language=JavaScript, caption={Menampilkan daftar dengan \texttt{map()}}, label={lst:13-komponen-map}]
function DaftarMahasiswa({ mahasiswa }) {
    return (
        <div className="bg-white rounded-lg shadow p-4">
            <h2 className="text-lg font-bold mb-2">
                Daftar Mahasiswa
            </h2>
            <ul>
                {mahasiswa.map((nama, index) => (
                    <li key={index} className="py-1 border-b">
                        {index + 1}. {nama}
                    </li>
                ))}
            </ul>
        </div>
    );
}

// Cara menggunakan:
// <DaftarMahasiswa
//     mahasiswa={["Budi", "Ani", "Andi", "Sari"]} />
\end{lstlisting}

\begin{catatan}
\texttt{map()} dalam JSX mirip \texttt{forEach()} di JavaScript biasa---tetapi mengembalikan nilai (\emph{return}). \texttt{key} diperlukan oleh React untuk mengidentifikasi setiap elemen dalam daftar.
\end{catatan}

\subsection{Kode 4 --- Komponen Bersarang (Nested)}

\begin{lstlisting}[language=JavaScript, caption={Komponen \texttt{Tombol} yang digunakan di dalam \texttt{HalamanUtama}}, label={lst:13-komponen-nested}]
// Komponen kecil -- tombol
function Tombol({ teks, warna }) {
    return (
        <button className={
            `px-4 py-2 rounded text-white font-bold ${warna}`
        }>
            {teks}
        </button>
    );
}

// Komponen besar -- menggunakan Tombol
function HalamanUtama() {
    return (
        <div className="p-6">
            <h1 className="text-2xl font-bold mb-4">
                Selamat Datang
            </h1>
            <p className="mb-4">Pilih aksi berikut:</p>
            {/* Menggunakan komponen Tombol */}
            <Tombol teks="Mulai" warna="bg-blue-500" />
            <span> </span>
            <Tombol teks="Batal" warna="bg-red-500" />
        </div>
    );
}
\end{lstlisting}

\begin{tip}
\texttt{HalamanUtama} menggunakan komponen \texttt{Tombol} dua kali dengan data berbeda. Ini seperti menggunakan fungsi dalam fungsi---komponen React bisa dipanggil berkali-kali dengan data berbeda.
\end{tip}

% =====================================================================
\section{JSX: HTML yang ``Bisa Dikode''}
\label{sec:13-jsx}
% =====================================================================

\textbf{JSX (JavaScript XML)} adalah sintaks yang mirip HTML tapi ditulis \textbf{di dalam JavaScript}. React menggunakan JSX sebagai cara mendeskripsikan tampilan UI.

\subsection{Perbedaan JSX vs HTML}

\begin{table}[htbp]
\centering
\caption{Perbedaan sintaks JSX dan HTML}
\label{tab:13-jsx-vs-html}
\begin{tabular}{lll}
\toprule
\textbf{Aspek} & \textbf{HTML} & \textbf{JSX} \\
\midrule
\emph{class} & \texttt{class="..."} & \texttt{className="..."} \\
\emph{for} (label) & \texttt{for="..."} & \texttt{htmlFor="..."} \\
\emph{style} & \texttt{style="color: red;"} & \texttt{style=\{\{ color: 'red' \}\}} \\
\emph{self-closing} & \texttt{<br>}, \texttt{<img>} & \texttt{<br />}, \texttt{<img />} \\
Ekspresi & Tidak bisa & Bisa pakai \texttt{\{\}} \\
Komentar & \texttt{<!-- ... -->} & \texttt{/* ... */} \\
\bottomrule
\end{tabular}
\end{table}

\subsection{JSX Adalah Ekspresi JavaScript}

\begin{lstlisting}[language=JavaScript, caption={Ekspresi JavaScript di dalam JSX}, label={lst:13-jsx-ekspresi}]
// Anda bisa menulis JavaScript di dalam JSX menggunakan {}

function Profil({ nama, umur }) {
    return (
        <div>
            <h2>Nama: {nama}</h2>
            <p>Umur: {umur} tahun</p>
            <p>Status: {umur >= 17 ? 'Dewasa' : 'Anak-anak'}</p>
            <p>Tahun Lahir: {2025 - umur}</p>
        </div>
    );
}

// {} di JSX = "masukkan hasil ekspresi JavaScript di sini"
\end{lstlisting}

\begin{catatan}
Perbandingan dengan JavaScript biasa: \texttt{"<h2>Nama: " + nama + "</h2>"} vs \texttt{<h2>Nama: \{nama\}</h2>}---JSX jauh lebih bersih dan mudah dibaca.
\end{catatan}

% =====================================================================
\section{Komponen: Blok Bangunan UI}
\label{sec:13-komponen}
% =====================================================================

\textbf{Analogi LEGO:} setiap komponen React adalah bata LEGO---punya bentuk dan warna tersendiri (\emph{props} \& \emph{style}), bisa disusun menjadi struktur lebih besar, dan bisa digunakan berulang kali di tempat berbeda.

\begin{lstlisting}[language=JavaScript, caption={Tiga pola komponen: tanpa props, dengan props, dan komposisi}, label={lst:13-pola-komponen}]
// 1. Komponen tanpa props (self-contained)
function Header() {
    return (
        <nav className="bg-blue-600 text-white p-4">
            <h1 className="text-xl font-bold">My Website</h1>
        </nav>
    );
}

// 2. Komponen dengan props
function Card({ title, description, buttonText }) {
    return (
        <div className="bg-white rounded-lg shadow p-4">
            <h2 className="text-lg font-bold">{title}</h2>
            <p className="text-gray-600 mb-3">{description}</p>
            <button className="bg-blue-500 text-white
                               px-4 py-2 rounded">
                {buttonText}
            </button>
        </div>
    );
}

// 3. Komponen menyusun komponen lain (Komposisi)
function App() {
    return (
        <div>
            <Header />
            <Card title="Tentang Kami"
                  description="Kami adalah..."
                  buttonText="Selengkapnya" />
            <Card title="Layanan"
                  description="Kami menyediakan..."
                  buttonText="Lihat" />
        </div>
    );
}
\end{lstlisting}

\begin{tip}
Perubahan di satu definisi komponen berlaku untuk \emph{semua} penggunaannya. Misal: jika Anda mengubah struktur \texttt{Card}, semua kartu otomatis berubah. Tanpa komponen, Anda harus mengubah setiap kartu secara manual.
\end{tip}

% =====================================================================
\section{Props: Mengirim Data ke Komponen}
\label{sec:13-props}
% =====================================================================

\textbf{Props (properties)} adalah cara mengirim data \textbf{dari luar ke dalam} komponen---seperti parameter fungsi.

\begin{lstlisting}[language=JavaScript, caption={Props sebagai ``parameter'' komponen}, label={lst:13-props-dasar}]
function Greeting({ nama }) {
    return <h1>Halo, {nama}!</h1>;
}

// Menggunakan komponen:
// <Greeting nama="Budi" />   -> render: <h1>Halo, Budi!</h1>
// <Greeting nama="Ani" />    -> render: <h1>Halo, Ani!</h1>
\end{lstlisting}

\subsection{Analogi Fungsi JavaScript}

\begin{lstlisting}[language=JavaScript, caption={Perbandingan fungsi JS biasa dengan komponen React}, label={lst:13-analogi}]
// JavaScript biasa
function sapa(nama) {
    return "Halo, " + nama + "!";
}

sapa("Budi");  // "Halo, Budi!"
sapa("Ani");   // "Halo, Ani!"

// React = konsep yang sama, tapi output-nya JSX (UI)
// <Greeting nama="Budi" /> = sapa("Budi") dalam konteks UI
\end{lstlisting}

\subsection{Props Bisa Berisi Apa Saja}

\begin{lstlisting}[language=JavaScript, caption={\emph{Props} dengan berbagai tipe data}, label={lst:13-props-lengkap}]
function ProductCard({ name, price, inStock, rating, image }) {
    return (
        <div className="bg-white rounded-lg shadow p-4">
            <img src={image} alt={name}
                 className="w-full h-48 object-cover rounded" />
            <h2 className="text-lg font-bold mt-2">{name}</h2>
            <p className="text-green-600 font-semibold">
                Rp {price.toLocaleString('id-ID')}
            </p>
            <p className="text-sm text-gray-500">
                {inStock ? 'Tersedia' : 'Habis'}
            </p>
            <p className="text-sm">{rating}/5</p>
        </div>
    );
}

// Menggunakan:
// <ProductCard name="Laptop ASUS" price={8500000}
//     inStock={true} rating={4.5} image="/laptop.jpg" />
\end{lstlisting}

\subsection{Aturan Penting Props}

\begin{lstlisting}[language=JavaScript, caption={Aturan \emph{props}: read-only, default, dan children}, label={lst:13-props-aturan}]
// 1. Props bersifat READ-ONLY -- komponen TIDAK BOLEH
//    mengubah props
function Avatar({ nama }) {
    // nama = "Hacker";  // JANGAN! Props tidak boleh diubah
    // Gunakan variabel baru jika perlu
    const namaBesar = nama.toUpperCase();
    return <div>{namaBesar}</div>;
}

// 2. Default props
function Button({ text = "Klik", color = "blue" }) {
    return <button className={`bg-${color}-500 text-white p-2`}>
        {text}
    </button>;
}

// 3. Children props -- konten di antara tag pembuka dan penutup
function Container({ children }) {
    return <div className="max-w-md mx-auto bg-white
                           p-6 rounded-lg shadow">
        {children}
    </div>;
}

// Menggunakan children:
// <Container>
//     <h1>Judul</h1>
//     <p>Paragraf di dalam container</p>
// </Container>
\end{lstlisting}

% =====================================================================
\section{Demo Lengkap: Membaca \& Memahami Kode React}
\label{sec:13-demo-lengkap}
% =====================================================================

Berikut kode React yang \textbf{lebih realistis}---seperti yang mungkin Anda temui di proyek nyata.

\begin{langkah}
  \item Tambahkan kode berikut ke \texttt{react-preview.html} (setelah kode-kode sebelumnya).
  \item Buka di peramban dan baca kode dengan seksama.
  \item Coba identifikasi pertanyaan-pertanyaan pemahaman di bawah kode.
\end{langkah}

\begin{lstlisting}[language=JavaScript, caption={Demo lengkap: Halaman Profil Mahasiswa dalam React}, label={lst:13-demo-lengkap}]
// =============================================
// File: App.jsx -- Halaman Profil Mahasiswa
// =============================================

// 1. Komponen kecil: badge status
function StatusBadge({ aktif }) {
    if (aktif) {
        return <span className="bg-green-100 text-green-800
                     px-3 py-1 rounded-full text-sm">
            Aktif
        </span>;
    }
    return <span className="bg-red-100 text-red-800
                 px-3 py-1 rounded-full text-sm">
        Tidak Aktif
    </span>;
}

// 2. Komponen medium: kartu profil
function KartuProfil({ nama, nim, jurusan, ipk, aktif }) {
    return (
        <div className="bg-white rounded-xl shadow-lg
                        p-6 max-w-sm">
            {/* Avatar */}
            <div className="w-20 h-20 rounded-full bg-blue-500
                flex items-center justify-center text-white
                text-2xl font-bold mx-auto mb-4">
                {nama.charAt(0)}
            </div>

            {/* Info */}
            <h2 className="text-xl font-bold text-center mb-1">
                {nama}
            </h2>
            <p className="text-gray-500 text-center mb-3">
                NIM: {nim}
            </p>
            <div className="text-center mb-3">
                <StatusBadge aktif={aktif} />
            </div>

            {/* Detail */}
            <div className="border-t pt-3 mt-3 space-y-1
                            text-sm">
                <p><strong>Jurusan:</strong> {jurusan}</p>
                <p><strong>IPK:</strong>
                    <span className={
                        ipk >= 3.5
                            ? "text-green-600"
                            : "text-red-600"
                    }>
                        {ipk}
                    </span>
                </p>
            </div>
        </div>
    );
}

// 3. Komponen utama: halaman
function HalamanProfil() {
    const daftarMahasiswa = [
        { nama: "Budi Santoso", nim: "2401001",
          jurusan: "Informatika", ipk: 3.8, aktif: true },
        { nama: "Ani Wijaya", nim: "2401002",
          jurusan: "Informatika", ipk: 3.2, aktif: true },
        { nama: "Andi Pratama", nim: "2401003",
          jurusan: "Sistem Informasi", ipk: 2.9, aktif: false },
    ];

    return (
        <div className="min-h-screen bg-gray-100 py-10">
            <h1 className="text-3xl font-bold text-center mb-8">
                Daftar Profil Mahasiswa
            </h1>
            <div className="flex flex-wrap justify-center
                            gap-6 px-4">
                {daftarMahasiswa.map((mhs) => (
                    <KartuProfil
                        key={mhs.nim}
                        nama={mhs.nama}
                        nim={mhs.nim}
                        jurusan={mhs.jurusan}
                        ipk={mhs.ipk}
                        aktif={mhs.aktif}
                    />
                ))}
            </div>
        </div>
    );
}
\end{lstlisting}

\begin{tangkapanlayar}
  {assets/images/minggu-13/react-kode-lengkap.png}
  {Kode lengkap di peramban dengan \emph{syntax highlighting}.}
  {Buka \texttt{react-preview.html} di peramban setelah menambahkan kode di atas.}
\end{tangkapanlayar}

\textbf{Checklist pemahaman}---centang jika sudah bisa menjawab:

\begin{itemize}
  \item[$\square$] Komponen \texttt{StatusBadge} menerima \emph{props} apa?
  \item[$\square$] Di mana \texttt{StatusBadge} digunakan (disebut)?
  \item[$\square$] Bagaimana cara menentukan warna IPK (hijau/merah)?
  \item[$\square$] Bagaimana \texttt{HalamanProfil} membuat daftar kartu?
  \item[$\square$] Apa fungsi \texttt{key} pada \texttt{KartuProfil}?
\end{itemize}

\begin{tangkapanlayar}
  {assets/images/minggu-13/react-checklist.png}
  {Checklist pemahaman yang sudah diisi.}
  {Tuliskan jawaban Anda untuk setiap pertanyaan di atas.}
\end{tangkapanlayar}

% =====================================================================
\section{Latihan}
\label{sec:13-latihan}
% =====================================================================

\begin{latihan}[Membaca Kode React]
Baca kode React berikut dan jawab pertanyaan di bawahnya:

\begin{lstlisting}[language=JavaScript, label={lst:13-latihan-kode}]
function TodoItem({ text, completed }) {
    return (
        <div className="flex items-center gap-2 p-2">
            <input type="checkbox" checked={completed}
                   readOnly />
            <span className={
                completed ? "line-through text-gray-400" : ""
            }>
                {text}
            </span>
        </div>
    );
}

function TodoList() {
    const todos = [
        { text: "Belajar HTML", completed: true },
        { text: "Belajar CSS", completed: true },
        { text: "Belajar JavaScript", completed: false },
        { text: "Belajar React", completed: false },
    ];

    return (
        <div className="max-w-md mx-auto bg-white p-4
                        rounded-lg shadow">
            <h2 className="text-xl font-bold mb-3">
                To-Do List
            </h2>
            {todos.map((todo, index) => (
                <TodoItem
                    key={index}
                    text={todo.text}
                    completed={todo.completed}
                />
            ))}
        </div>
    );
}
\end{lstlisting}

Pertanyaan:
\begin{enumerate}
  \item Komponen \texttt{TodoItem} menerima \emph{props} apa saja?
  \item Apa yang terjadi pada teks jika \texttt{completed === true}?
  \item Berapa banyak \texttt{TodoItem} yang akan ditampilkan?
  \item Berapa banyak item yang sudah selesai (centang)?
  \item Jika Anda ingin menambah item baru, bagian mana yang perlu diubah?
\end{enumerate}

\begin{tangkapanlayar}
  {assets/images/minggu-13/latihan1-jawaban.png}
  {Tuliskan jawaban Anda (boleh di file teks atau di komentar di dalam file HTML).}
  {Jawab setiap pertanyaan di atas berdasarkan kode yang diberikan.}
\end{tangkapanlayar}
\end{latihan}

\begin{latihan}[Menulis Pseudo-Code React]
Bayangkan Anda diminta membuat komponen \texttt{KartuProduk} untuk toko \emph{online}. \textbf{Tanpa menjalankan kode}, tuliskan (dalam bentuk \emph{pseudo-code} atau kode JSX mentah) komponen yang:
\begin{enumerate}
  \item Menerima \emph{props}: \texttt{nama}, \texttt{harga}, \texttt{gambar} (URL), \texttt{diskon} (persen, bisa 0).
  \item Menampilkan gambar, nama produk, harga asli, dan harga setelah diskon (jika diskon $>$ 0).
  \item Menggunakan Tailwind CSS untuk \emph{styling} (minimal: \texttt{rounded}, \texttt{shadow}, \texttt{padding}).
\end{enumerate}

\textbf{Output yang diharapkan (kira-kira tampilan yang Anda bayangkan):}

\begin{lstlisting}
+------------------------+
|   [Gambar Produk]      |
|                        |
|   Laptop ASUS          |
|   Rp 8.500.000         |
|   ~~Rp 10.000.000~~    |  <-- hanya jika diskon > 0
|   (hemat 15%)          |
+------------------------+
\end{lstlisting}

\begin{tangkapanlayar}
  {assets/images/minggu-13/latihan2-pseudocode.png}
  {Kode JSX \emph{pseudo-code} yang Anda tulis.}
  {Tuliskan kode JSX untuk komponen \texttt{KartuProduk} Anda.}
\end{tangkapanlayar}
\end{latihan}

\begin{latihan}[Peta Pemahaman --- Opsional]
Buat daftar 5 konsep React yang sudah Anda pahami dan 5 konsep yang masih ingin dipelajari lebih lanjut. Ini akan menjadi acuan saat Anda belajar React di semester lanjutan.

\begin{center}
\textbf{Tabel: Contoh format peta pemahaman}
\medskip
\begin{tabular}{ll}
\toprule
\textbf{Sudah Dipahami} & \textbf{Ingin Dipelajari} \\
\midrule
Komponen adalah fungsi & \texttt{useState} hook \\
\emph{Props} = parameter komponen & \texttt{useEffect} (\emph{side effects}) \\
JSX mirip HTML tapi ada bedanya & React Router (navigasi) \\
Komponen bisa bersarang & API \emph{calls} di React \\
\texttt{map()} untuk membuat daftar & Context API (\emph{state} global) \\
\bottomrule
\end{tabular}
\end{center}

\begin{tangkapanlayar}
  {assets/images/minggu-13/latihan3-peta.png}
  {Daftar pemahaman Anda.}
  {Tuliskan 5 konsep yang sudah dipahami dan 5 yang ingin dipelajari.}
\end{tangkapanlayar}
\end{latihan}

% =====================================================================
\section{Troubleshooting \& Catatan}
\label{sec:13-troubleshooting}
% =====================================================================

\subsection{FAQ Umum tentang React}

\begin{catatan}
\textbf{Q: Apakah saya harus menguasai React sekarang?}
Tidak. Modul ini hanya pengenalan. Anda akan belajar React secara mendalam di semester lanjutan. Yang penting sekarang adalah \textbf{membaca dan memahami konsep dasarnya}.
\end{catatan}

\begin{catatan}
\textbf{Q: Kenapa ada banyak simbol aneh di kode React?}
\begin{itemize}
  \item \texttt{\{\}} = ekspresi JavaScript di dalam JSX
  \item \texttt{() => \{\}} = \emph{arrow function} (sudah dipelajari di Minggu 10)
  \item \texttt{className} = \emph{attribute} di JSX untuk menetapkan class CSS (karena \texttt{class} adalah kata kunci JavaScript)
\end{itemize}
\end{catatan}

\begin{catatan}
\textbf{Q: React vs Vue vs Angular---mana yang terbaik?}
Tidak ada yang ``terbaik'' secara mutlak. React paling banyak digunakan di industri global. Vue lebih populer di beberapa region. Angular untuk aplikasi \emph{enterprise}. Untuk kuliah, React adalah pilihan aman karena ekosistemnya paling besar.
\end{catatan}

\subsection{Kosa Kata React}

\begin{table}[htbp]
\centering
\caption{Istilah penting dalam React}
\label{tab:13-kosakata}
\begin{tabular}{ll}
\toprule
\textbf{Istilah} & \textbf{Arti} \\
\midrule
\emph{Component} & Blok bangunan UI---fungsi yang mengembalikan JSX \\
JSX & JavaScript XML---sintaks HTML-\emph{like} di dalam JavaScript \\
\emph{Props} & Data yang dikirim ke komponen (\emph{read-only}) \\
\emph{State} & Data internal komponen yang bisa berubah (semester lanjutan) \\
\emph{Virtual DOM} & Representasi DOM di memori---React update DOM secara efisien \\
\emph{Hook} & Fitur khusus React (\texttt{useState}, \texttt{useEffect})---semester lanjutan \\
\bottomrule
\end{tabular}
\end{table}

\subsection{Ringkasan Perjalanan Minggu 10--13}

\begin{table}[htbp]
\centering
\caption{Topik dan keterampilan yang diperoleh dari Minggu 10--13}
\label{tab:13-ringkasan}
\begin{tabular}{lll}
\toprule
\textbf{Minggu} & \textbf{Topik} & \textbf{Skill yang Didapat} \\
\midrule
10 & JavaScript dasar & Variabel, fungsi, kondisi, DOM selection \\
11 & JavaScript interaktif & \emph{Event handling}, manipulasi DOM \\
12 & Tailwind CSS & \emph{Utility-first CSS}, layout \emph{responsive} \\
13 & React \emph{preview} & Membaca kode komponen, JSX, \emph{props} \\
\bottomrule
\end{tabular}
\end{table}

\begin{tip}
\textbf{Next step (semester lanjutan):} Anda akan membangun aplikasi React yang sebenarnya---menggunakan \emph{state}, \emph{effect}, \emph{routing}, dan API. Fondasi dari Minggu 10--13 akan membuat proses itu jauh lebih mudah!
\end{tip}
